\chapter{Distribusi Binomial}\label{ch:g11:binom}

Ulangilah percobaan ya/tidak yang sama beberapa kali, secara saling bebas, lalu
bilanglah keberhasilannya: \omterm{def:g11:prob:rv}{distribusi} yang muncul --- yaitu \emph{binomial} ---
adalah \omterm{def:g11:prob:rv}{distribusi} diskret yang paling penting di antara semuanya. Bab ini
membangunnya dengan \omterm{met:g10:proba:tree}{pohon} dan pembilangan jalur; rumus tertutup bagi cacah
jalurnya (dengan faktorial) datang bersama perkakas pembilangan pada
\cref{ch:g12:comb}, dan \omterm{def:g11:prob:rv}{distribusinya} ditinjau ulang pada \cref{ch:g12:randvar}.

\section{Percobaan Bernoulli}

\begin{definition}[Percobaan Bernoulli]\label{def:g11:binom:bernoulli}
Sebuah \emph{percobaan Bernoulli}\index{percobaan Bernoulli} adalah percobaan
dengan tepat dua hasil: \emph{berhasil}, dengan peluang $p$, dan
\emph{gagal}, dengan peluang $1 - p$. \omterm{def:g11:prob:rv}{Peubah acak} $X$ yang bernilai
$1$ pada keberhasilan dan $0$ pada kegagalan dikatakan mengikuti
\emph{\omterm{def:g11:prob:rv}{distribusi} Bernoulli} $\mathcal B(p)$; lalu
\[
\E(X) = p,
\qquad
\V(X) = p(1 - p).
\]
\end{definition}

\begin{proof}[Bukti kedua rumusnya]
$\E(X) = p \times 1 + (1-p) \times 0 = p$; dan karena $X^2 = X$ (baik $0$
maupun $1$ adalah kuadrat dirinya sendiri), $\E(X^2) = p$, sehingga menurut
\cref{prop:g11:prob:konig}, $\V(X) = p - p^2 = p(1-p)$.
\end{proof}

\begin{definition}[Percobaan berulang yang saling bebas]\label{def:g11:binom:repeated}
Mengulang sebuah \omterm{def:g11:binom:bernoulli}{percobaan Bernoulli} sebanyak $n$ kali \emph{secara saling
bebas} berarti: hasil setiap percobaannya tidak berpengaruh pada yang lain, dan
peluang sebarang runtunan hasil yang lengkap adalah \emph{hasil kali}
peluang sepanjang jalur \omterm{met:g10:proba:tree}{pohon} yang bersesuaian --- $p$ untuk setiap
keberhasilan, $1 - p$ untuk setiap kegagalan.
\end{definition}

\begin{example}\label{ex:g11:binom:threetrials}
Tiga percobaan saling bebas dengan peluang berhasil $p$. Runtunan BGB
(berhasil, gagal, berhasil) berpeluang $p(1-p)p = p^2(1-p)$ --- dan
demikian pula \emph{setiap} runtunan dengan tepat dua keberhasilan, tanpa
memandang letaknya: hanya \emph{banyaknya} B dan G yang berpengaruh.
\end{example}

\section{Cacah jalur dan koefisien binomial}

\begin{definition}[Koefisien binomial]\label{def:g11:binom:coefficient}
Pada \omterm{met:g10:proba:tree}{pohon} $n$ percobaan saling bebas, \emph{koefisien
binomial}\index{koefisien binomial} $\binom{n}{k}$ (dibaca ``$n$
pilih $k$'') adalah banyaknya jalur yang memuat tepat $k$ keberhasilan.
\end{definition}

\begin{example}
$\binom{3}{2} = 3$: yaitu jalur BBG, BGB, GBB. Demikian pula $\binom{3}{0} = 1$
(jalur GGG), $\binom{3}{1} = 3$ dan $\binom{3}{3} = 1$. Menurut
kelaziman dan menurut \omterm{met:g10:proba:tree}{pohonnya}, $\binom n0 = \binom nn = 1$ untuk setiap $n$.
\end{example}

\begin{omfigure}
\begin{tikzpicture}[font=\small]
  \coordinate (root) at (0,0);
  \node (S) at (1.5,1.5) {B};
  \node (F) at (1.5,-1.5) {G};
  \node (SS) at (3.0,2.25) {B};
  \node (SF) at (3.0,0.75) {G};
  \node (FS) at (3.0,-0.75) {B};
  \node (FF) at (3.0,-2.25) {G};
  \node (SSS) at (4.5,2.6) {B};
  \node (SSF) at (4.5,1.85) {G};
  \node (SFS) at (4.5,1.1) {B};
  \node (SFF) at (4.5,0.35) {G};
  \node (FSS) at (4.5,-0.35) {B};
  \node (FSF) at (4.5,-1.1) {G};
  \node (FFS) at (4.5,-1.85) {B};
  \node (FFF) at (4.5,-2.6) {G};
  \draw[omThm, thick] (root) -- (S);
  \draw[omThm, thick] (S) -- (SS);   \draw[omThm, thick] (SS) -- (SSF);
  \draw[omThm, thick] (S) -- (SF);   \draw[omThm, thick] (SF) -- (SFS);
  \draw[omThm, thick] (root) -- (F);
  \draw[omThm, thick] (F) -- (FS);   \draw[omThm, thick] (FS) -- (FSS);
  \draw[gray] (SS) -- (SSS);
  \draw[gray] (SF) -- (SFF);
  \draw[gray] (FS) -- (FSF);
  \draw[gray] (F) -- (FF);
  \draw[gray] (FF) -- (FFS);
  \draw[gray] (FF) -- (FFF);
  \node[gray, anchor=west] at (5.0,2.6) {$3$ keberhasilan};
  \node[omThm, anchor=west] at (5.0,1.85) {$2$ keberhasilan};
  \node[omThm, anchor=west] at (5.0,1.1) {$2$ keberhasilan};
  \node[gray, anchor=west] at (5.0,0.35) {$1$ keberhasilan};
  \node[omThm, anchor=west] at (5.0,-0.35) {$2$ keberhasilan};
  \node[gray, anchor=west] at (5.0,-1.1) {$1$ keberhasilan};
  \node[gray, anchor=west] at (5.0,-1.85) {$1$ keberhasilan};
  \node[gray, anchor=west] at (5.0,-2.6) {$0$ keberhasilan};
\end{tikzpicture}

{\small \omterm{met:g10:proba:tree}{Pohon} $n = 3$ percobaan: $\binom{3}{2} = 3$ jalur (merah) membawa
tepat dua keberhasilan, masing-masing berpeluang $p^2(1-p)$.}
\end{omfigure}

\begin{proposition}[Aturan Pascal]\label{prop:g11:binom:pascal}
\index{segitiga Pascal}
Untuk $1 \leq k \leq n - 1$:
\[
\binom{n}{k} = \binom{n-1}{k-1} + \binom{n-1}{k}.
\]
\end{proposition}

\begin{proof}
Kelompokkan jalur pada \omterm{met:g10:proba:tree}{pohon} $n$ percobaan yang mempunyai $k$ keberhasilan
menurut percobaan \emph{terakhirnya}. Jalur yang berakhir pada keberhasilan
diperoleh dari jalur $n-1$ percobaan pertama dengan $k - 1$ keberhasilan: ada
$\binom{n-1}{k-1}$ jalur semacam itu. Jalur yang berakhir pada kegagalan
memperpanjang jalur dengan $k$ keberhasilan di antara $n-1$ percobaan pertama:
ada $\binom{n-1}{k}$ jalur. Setiap jalurnya termasuk tepat satu di antara kedua jenis itu.
\end{proof}

Aturan Pascal membangkitkan koefisiennya baris demi baris --- setiap entrinya
adalah jumlah dua entri di atasnya:
\[
\begin{array}{ccccccccccc}
&&&&&1&&&&&\\
&&&&1&&1&&&&\\
&&&1&&2&&1&&&\\
&&1&&3&&3&&1&&\\
&1&&4&&6&&4&&1&\\
1&&5&&10&&10&&5&&1
\end{array}
\]

\begin{remark}
Rumus tertutupnya, $\binom nk = \frac{n!}{k!(n-k)!}$, beserta teori
pembilangan yang sistematis, ditetapkan pada \cref{ch:g12:comb}. Pada
tingkat ini, \omterm{prop:g11:binom:pascal}{segitiga Pascal} menghitung setiap koefisien yang kita perlukan.
\end{remark}

\section{Distribusi binomial}

\begin{theorem}[Distribusi binomial]\label{thm:g11:binom:binomial}
Misalkan $X$ membilang keberhasilan pada $n$ \omterm{def:g11:binom:bernoulli}{percobaan Bernoulli} saling bebas
berparameter $p$. Maka $X$ mengikuti \emph{distribusi
binomial}\index{distribusi binomial} $\mathcal B(n, p)$:
\[
\P(X = k) = \binom{n}{k}\, p^k (1-p)^{n-k},
\qquad k = 0, 1, \dots, n .
\]
\end{theorem}

\begin{proof}
\omterm{def:g11:prob:model}{Kejadian} $X = k$ adalah kumpulan semua jalur dengan tepat $k$
keberhasilan. Setiap jalur semacam itu berpeluang $p^k(1-p)^{n-k}$: hasil kali
sepanjang jalurnya memuat $k$ faktor $p$ dan $n - k$ faktor $1-p$, dalam
urutan tertentu (\cref{def:g11:binom:repeated}). Ada $\binom nk$ jalur semacam
itu (\cref{def:g11:binom:coefficient}), dan peluangnya dijumlahkan.
\end{proof}

\begin{example}\label{ex:g11:binom:quiz}
Sebuah kuis mempunyai $5$ soal saling bebas, masing-masing dengan $4$ pilihan;
seorang siswa menjawab secara acak, sehingga setiap soal berhasil dengan
$p = \frac14$. Banyaknya jawaban benar $X$ mengikuti
$\mathcal B\left(5, \frac14\right)$, dan, dengan baris $5$ \omterm{prop:g11:binom:pascal}{segitiga Pascal}:
\[
\P(X = 2) = \binom52 \left(\frac14\right)^2\left(\frac34\right)^3
= 10 \times \frac{1}{16} \times \frac{27}{64}
= \frac{270}{1024} \approx 0.26 .
\]
Peluang memperoleh paling sedikit satu jawaban benar dihitung dengan \omterm{def:g10:proba:operations}{komplemennya}:
$\P(X \geq 1) = 1 - \P(X = 0) = 1 - \left(\frac34\right)^5 \approx
0.76$.
\end{example}

\begin{proposition}[Nilai harapan dan ragam]\label{prop:g11:binom:expectation}
Jika $X \sim \mathcal B(n, p)$:
\[
\E(X) = np,
\qquad
\V(X) = np(1-p).
\]
\end{proposition}

\begin{proof}[Penjelasan]
Tulislah $X = X_1 + X_2 + \dots + X_n$, dengan $X_i$ bernilai $1$ bila
percobaan ke-$i$ berhasil: setiap $X_i$ adalah peubah Bernoulli dengan \omterm{def:g11:prob:expectation}{nilai
harapan} $p$ (\cref{def:g11:binom:bernoulli}). \omterm{def:g11:stat:mean}{Rata-rata} dapat dijumlahkan ---
menjumlahkan $n$ sumbangannya memberi $\E(X) = np$. Bahwa \emph{\omterm{def:g11:prob:variance}{ragam}} juga
dapat dijumlahkan untuk peubah yang saling bebas memang benar tetapi lebih
peka: rumus \omterm{def:g11:prob:variance}{ragamnya} \emph{diterima tanpa bukti pada tingkat ini} dan
dibuktikan pada \cref{ch:g12:sums}.
\end{proof}

\begin{omfigure}
\begin{tikzpicture}
\begin{axis}[
  width=10cm, height=5cm,
  ybar, bar width=9pt,
  xlabel={$k$}, ylabel={$\P(X = k)$},
  xmin=-0.8, xmax=10.8, ymin=0, ymax=0.3,
  xtick={0,1,2,3,4,5,6,7,8,9,10},
  ytick={0.1, 0.2},
  tick label style={font=\small},
  label style={font=\small},
  title={$\mathcal B(10,\ 0.5)$}]
  \addplot[fill=omDef!40, draw=omDef] coordinates {
    (0,0.00098) (1,0.00977) (2,0.04395) (3,0.11719) (4,0.20508)
    (5,0.24609) (6,0.20508) (7,0.11719) (8,0.04395) (9,0.00977)
    (10,0.00098)};
\end{axis}
\end{tikzpicture}

{\small \omterm{def:g11:prob:rv}{Distribusi} $\mathcal B(10, 0.5)$ (sepuluh lemparan koin setimbang):
berpusat di $\E(X) = np = 5$, setangkup, dengan hampir seluruh peluangnya
berada antara $2$ dan $8$.}
\end{omfigure}

\begin{method}[Mengenali situasi binomial]\label{met:g11:binom:recognize}
Sebelum menulis $X \sim \mathcal B(n, p)$, periksalah tiga bahannya: sebuah
\emph{banyaknya percobaan yang tetap} $n$, yang ditetapkan lebih dahulu; setiap
percobaan mempunyai \emph{dua hasil} dengan peluang berhasil $p$ yang
\emph{sama}; dan percobaannya \emph{saling bebas} (dengan pengembalian, atau
dengan alat yang terpisah). Pengambilan tanpa pengembalian dari populasi kecil
\emph{bukan} binomial --- peluangnya berubah pada setiap pengambilan
(\cref{exo:g11:prob:6}).
\end{method}

\section{Penyampelan: apakah pengamatannya mengejutkan?}

\omterm{def:g11:prob:rv}{Distribusi} binomial menjawab sebuah pertanyaan yang sangat praktis:
\emph{jika peluang berhasilnya sungguh $p$, cacah keberhasilan mana saja yang
masuk akal?}

\begin{example}\label{ex:g11:binom:decision}
Sebuah mesin seharusnya menghasilkan paling banyak $10\%$ barang cacat. Pada
sebuah kumpulan berisi $10$ barang, $4$ di antaranya cacat. Nasib buruk atau
mesin rusak? Jika mesinnya baik-baik saja, banyaknya barang cacat mengikuti
$\mathcal B(10, 0.1)$, dan
\[
\P(X \geq 4) = 1 - \P(X \leq 3) \approx 1 - 0.987 = 0.013 :
\]
yaitu sekitar satu peluang dari $80$. Mengamati \omterm{def:g11:prob:model}{kejadian} setidak mungkin ini
adalah isyarat yang kuat --- kita \emph{menolak} hipotesis bahwa mesinnya masih
bekerja pada $10\%$, sambil mengingat bahwa keputusan itu dapat keliru dengan
peluang sekitar $0.013$.
\end{example}

\begin{method}[Aturan keputusan dari model binomial]\label{met:g11:binom:decision}
Untuk menilai cacah keberhasilan $k$ yang teramati terhadap hipotesis
$X \sim \mathcal B(n, p)$: hitunglah peluang, di bawah hipotesis itu,
munculnya hasil yang \emph{paling sedikit seekstrem} $k$. Jika peluang itu
sangat kecil (kelaziman yang umum: di bawah $5\%$), tolaklah hipotesisnya;
bila tidak, pengamatannya sesuai dengan hipotesis itu. Ambangnya adalah sebuah
pilihan, bukan sebuah teorema --- statistika mengukur risikonya, dan pemakainya
menerima risiko itu.
\end{method}

\section{Latihan}

\begin{exercise}[$\star$]\label{exo:g11:binom:1}
Perluaslah \omterm{prop:g11:binom:pascal}{segitiga Pascal} sampai baris $7$, lalu berikan nilai
$\binom62$, $\binom{7}{3}$ dan $\binom74$.
\end{exercise}

\begin{exercise}[$\star$]\label{exo:g11:binom:2}
Sebuah dadu setimbang dilempar $4$ kali; $X$ membilang munculnya enam. Berikan
alasan bahwa $X \sim \mathcal B\left(4, \frac16\right)$ lalu hitunglah
$\P(X = 0)$, $\P(X = 1)$ dan $\P(X \geq 2)$.
\end{exercise}

\begin{exercise}[$\star$]\label{exo:g11:binom:3}
Manakah di antara berikut yang binomial? Berikan alasannya.
\begin{enumerate}
  \item Banyaknya gambar pada $20$ lemparan koin setimbang.
  \item Banyaknya kartu hati pada $5$ kartu yang dibagikan dari satu set kartu.
  \item Banyaknya hari hujan pekan depan, bila setiap harinya hujan dengan
        peluang $0.3$ secara saling bebas.
\end{enumerate}
\end{exercise}

\begin{exercise}[$\star$]\label{exo:g11:binom:4}
$X \sim \mathcal B(50, 0.2)$. Berikan $\E(X)$, $\V(X)$ dan $\sigma(X)$.
\end{exercise}

\begin{exercise}[$\star\star$]\label{exo:g11:binom:5}
Seorang pemanah mengenai sasaran dengan peluang $0.7$ pada setiap tembakan,
secara saling bebas. Pada $6$ tembakan, hitunglah peluang tepat $4$ kali kena,
dan peluang paling sedikit $5$ kali kena.
\end{exercise}

\begin{exercise}[$\star\star$]\label{exo:g11:binom:6}
Sebuah tes benar/salah mempunyai $8$ soal; seorang siswa menerka setiap
jawabannya. Berapa peluang ia lulus (paling sedikit $6$ jawaban benar)?
\end{exercise}

\begin{exercise}[$\star\star$]\label{exo:g11:binom:7}
Setiap kotak sereal yang dibeli berisi figur A atau figur B, masing-masing
dengan peluang $\frac12$, secara saling bebas. Seorang pengumpul membeli $5$
kotak. Hitunglah peluang bahwa pengumpul itu memperoleh paling sedikit satu
figur dari setiap jenis. (\omterm{def:g10:proba:operations}{Komplemennya}: semuanya A atau semuanya B.)
\end{exercise}

\begin{exercise}[$\star\star$]\label{exo:g11:binom:8}
Seorang pebasket memasukkan lemparan bebas dengan peluang $p$,
secara saling bebas. Misalkan $X \sim \mathcal B(3, p)$ banyaknya lemparan yang
masuk dari tiga lemparan. Nyatakan $\P(X = 3)$ dan $\P(X \geq 1)$ sebagai \omterm{def:g11:func:function}{fungsi}
$p$, lalu carilah untuk $p$ yang mana peluang ketiganya masuk sama dengan
$\frac{27}{64}$.
\end{exercise}

\begin{exercise}[$\star\star$]\label{exo:g11:binom:9}
Berapa kali sebuah koin setimbang harus dilempar agar peluang memperoleh
paling sedikit satu gambar melampaui $0.99$? (Pakailah \omterm{def:g10:proba:operations}{komplemennya}, lalu
cobalah nilai $n$ berturut-turut.)
\end{exercise}

\begin{exercise}[$\star\star$]\label{exo:g11:binom:10}
Dengan aturan Pascal (\cref{prop:g11:binom:pascal}) dan
$\binom n0 = \binom nn = 1$, buktikan bahwa entri setiap baris
\omterm{prop:g11:binom:pascal}{segitiga Pascal} berjumlah $2^n$: tafsirkan kedua ruasnya sebagai membilang
semua jalur pada \omterm{met:g10:proba:tree}{pohonnya}.
\end{exercise}

\begin{exercise}[$\star\star\star$]\label{exo:g11:binom:11}
Seorang politikus mengaku memperoleh dukungan $60\%$. Pada \omterm{def:g10:proba:sample}{sampel} acak berisi
$10$ orang, hanya $3$ yang mendukung.
\begin{enumerate}
  \item Di bawah pengakuan itu, \omterm{def:g11:prob:rv}{distribusi} apa yang diikuti banyaknya
        pendukung $X$ pada \omterm{def:g10:proba:sample}{sampelnya}? Hitunglah $\P(X \leq 3)$.
  \item Dengan aturan keputusan \cref{met:g11:binom:decision} pada ambang
        $5\%$, apakah pengamatannya sesuai dengan pengakuan itu?
\end{enumerate}
\end{exercise}

\section{Soal: Papan Galton}

\begin{problem}[{Soal akhir pekan --- bola, pasak dan segitiga
Pascal: bagaimana bentuk lonceng itu lahir, mengapa rangkaian laga pamungkas
menguntungkan tim yang lebih kuat, dan kapan harus berteriak curang}]
\label{pb:g11:binom:1}
Jatuhkanlah seribu bola melalui kisi pasak, setiap pantulannya sebuah
lemparan koin kiri-atau-kanan yang setimbang, dan kotak di bawahnya terisi
menjadi lonceng yang mulus dan setangkup --- setiap kali. Mesin itu disebut
papan Galton, dan matematikanya tepat sama dengan \omterm{def:g11:prob:rv}{distribusi} binomial bab ini
(\cref{thm:g11:binom:binomial}). Soal ini
membangun segitiganya, menjalankan papannya, mewasiti rangkaian
tujuh laga, lalu berakhir di tempat binomialnya memperoleh
upah: memutuskan kapan sebuah pengamatan seharusnya membuat kita meragukan
sebuah pengakuan.

\medskip
\noindent\textbf{Bagian I --- Segitiganya.}
\begin{enumerate}
  \item Bangunlah \omterm{prop:g11:binom:pascal}{segitiga Pascal} sampai baris $6$
        (\cref{prop:g11:binom:pascal}). Nyatakan lalu jelaskan
        kesimetrian $\binom nk = \binom{n}{n-k}$ dalam satu kalimat
        (memilih $k$ objek sama saja dengan\,\dots).
  \item Periksalah pada baris $4$ dan $5$ bahwa setiap barisnya berjumlah
        $2^n$, lalu buktikan: apa yang dibilang oleh semua $\binom nk$
        bersama-sama?
  \item Turunkan ulang aturan Pascal
        $\binom{n+1}{k} = \binom nk + \binom{n}{k-1}$ lewat
        alasan panitia: tetapkan satu orang istimewa lalu
        belahlah panitianya menurut nasib orang itu.
  \item Hitunglah $\binom73$ dua kali: dari segitiganya, dan dari
        rumus faktorialnya.
  \item Periksalah identitas tangga
        $\binom22 + \binom32 + \binom42 + \binom52 =
        \binom63$, lalu jelaskan lewat penerapan berjenjang aturan Pascal
        dari $\binom63$ ke bawah.
\end{enumerate}

\medskip
\noindent\textbf{Bagian II --- Papannya.}
Sebuah bola jatuh melalui $n$ baris pasak; pada setiap pasak bola itu memantul
ke kiri atau ke kanan dengan peluang $\frac12$, secara saling bebas. Nomorilah
kotaknya $0$ sampai $n$ menurut cacah pantulan ke kanan.
\begin{enumerate}[resume]
  \item Jelaskan, dengan daftar periksa
        \cref{met:g11:binom:recognize}, mengapa nomor kotaknya
        mengikuti \omterm{def:g11:prob:rv}{distribusi} binomial
        $\mathcal B\!\left(n, \frac12\right)$.
  \item Untuk papan kecil ($n = 4$): berikan kelima peluang
        kotaknya. Kotak mana yang paling padat?
  \item Kini $n = 10$ dan $1\,024$ bola: berapa \emph{harapan}
        cacah bola pada kotak tengah, pada kotak
        $7$, dan pada setiap kotak tepinya? Perikan bentuk tumpukannya.
  \item Untuk $X \sim \mathcal B\!\left(10, \frac12\right)$:
        hitunglah $\E(X)$, $V(X)$ dan $\sigma$
        (\cref{prop:g11:binom:expectation}); lalu hitunglah
        bagian bola yang diharapkan berada dalam $2\sigma$ dari
        pusatnya (kotak $2$ sampai $8$) lalu bandingkan dengan jaminan
        Chebyshev pada \cref{pb:g11:stat:1}.
  \item Sebuah papan yang dimiringkan memantul ke kanan dengan peluang $0.6$:
        berikan $\E$, $V$ dan $\sigma$ untuk $n = 10$, lalu perikan
        apa yang terjadi pada tumpukannya.
  \item Dalam satu atau dua kalimat: apa, dalam rancangan papan itu,
        yang membuat bentuk loncengnya --- dan mengapa begitu banyak
        besaran di dunia nyata (tinggi badan, galat pengukuran)
        menumpuk dengan cara yang sama? (Teorema mendalam di balik keduanya
        adalah teorema limit pusat, puncak kuliah peluang pada jilid
        universitas.)
\end{enumerate}

\medskip
\noindent\textbf{Bagian III --- Tujuh laga pamungkas.}
Dua tim memainkan sebuah rangkaian: yang lebih dahulu mencapai $4$ kemenangan
merebut gelarnya; setiap laganya saling bebas.
\begin{enumerate}[resume]
  \item Tim yang sepadan ($p = \frac12$): hitunglah peluang
        bahwa rangkaiannya berakhir dengan sapu bersih (tepat $4$ laga).
  \item Hitunglah peluang bahwa rangkaiannya berlangsung penuh
        $7$ laga (kedudukannya setelah $6$ laga harus berapa?).
  \item Lengkapkan \omterm{def:g11:prob:rv}{distribusi} panjang rangkaiannya
        ($4$, $5$, $6$ atau $7$ laga) untuk tim yang sepadan, lalu
        hitunglah harapan panjangnya. Panjang mana yang paling mungkin?
  \item Kini satu tim memenangkan setiap laga dengan $p = 0.6$. Hitunglah
        peluangnya merebut rangkaian itu (menang dalam $4$, $5$, $6$
        atau $7$ laga: pada setiap kasusnya tim itu memenangkan laga terakhir
        dan $3$ dari laga sebelumnya). Apa yang dilakukan rangkaian itu
        terhadap keunggulan per laganya?
  \item Bandingkan dengan satu laga final tunggal ($60\,\%$) dan sistem
        tiga laga (hitunglah). Nyatakan pengaruh umum panjang rangkaian
        terhadap keterampilan melawan keberuntungan --- dan mengapa liga
        lebih menyukai final yang panjang.
\end{enumerate}

\medskip
\noindent\textbf{Bagian IV --- Kapan harus berteriak curang.}
\begin{enumerate}[resume]
  \item Sebuah koin dilempar $100$ kali dan menunjukkan $62$ gambar. Untuk
        koin yang setimbang, berikan $\E$, $\sigma$, dan \omterm{pb:g11:stat:1}{skor-z}
        (\cref{pb:g11:stat:1}) pengamatannya. Bagaimana putusannya
        di bawah kelaziman $2\sigma$?
  \item Seorang pemasok mengaku bahwa $2\,\%$ suku cadangnya cacat. Pada
        sebuah kumpulan berisi $50$ kamu menemukan $3$ yang cacat. Hitunglah
        $\P(X \geq 3)$ di bawah pengakuan itu
        ($X \sim \mathcal B(50, 0.02)$; lewatilah
        $\P(X = 0), \P(X = 1), \P(X = 2)$). Mengkhawatirkan pada
        ambang $5\,\%$ (\cref{met:g11:binom:decision},
        \cref{exo:g11:binom:11})?
  \item Kegigihan pengundi: setiap karcis memenangkan (sesuatu) dengan
        peluang $\frac{1}{1000}$. Hitunglah peluang
        paling sedikit satu kemenangan dengan $1\,000$ karcis. Jawabannya
        ($\approx 63\,\%$, bukan $100\,\%$!) menyembunyikan sebuah tetapan
        termasyhur: hitunglah $0.999^{1000}$ lalu simpanlah bilangan
        $0.368$ dalam ingatan untuk tahun berikutnya.
  \item Penutup --- potret binomialnya: daftar periksa pengenalannya
        ($n$ yang tetap, kebebasan, $p$ yang tetap);
        \omterm{prop:g11:binom:pascal}{segitiga Pascal} sebagai tabelnya; lonceng sebagai bentuknya;
        $np$ dan $np(1-p)$ sebagai kompasnya; serta kedua pewarisnya yang
        menanti pada tahun berikutnya --- kurva lonceng yang mulus dan
        hukum bilangan besar. Masing-masing satu kalimat.

\end{enumerate}
\end{problem}
